\section*{Capítulo \ref{ch:g3:measure} --- Medidas}

\begin{solution}{exo:g3:measure:1}
Largura de uma unha: mm. Altura de um colega: cm (ou m). Pátio da
escola: m. Distância entre cidades: km.
\end{solution}

\begin{solution}{exo:g3:measure:2}
$5$ m $= 500$ cm; \quad $70$ mm $= 7$ cm; \quad
$4$ km $= 4\,000$ m; \quad $600$ cm $= 6$ m.
\end{solution}

\begin{solution}{exo:g3:measure:3}
As respostas dependem do triângulo; o comprimento total do contorno
é a \omterm{def:g1:addition:def}{soma} dos três lados medidos (esse total chama-se
\emph{perímetro}, estudado no \cref{ch:g4:measure}).
\end{solution}

\begin{solution}{exo:g3:measure:4}
$2$ m $+ 45$ cm $= 200 + 45 = 245$ cm.

$1$ km $+ 250$ m $= 1\,000 + 250 = 1\,250$ m.

$3$ cm $+ 8$ mm $= 30 + 8 = 38$ mm.
\end{solution}

\begin{solution}{exo:g3:measure:5}
$3$ kg $= 3\,000$ g; \quad $2\,500$ g $= 2$ kg e $500$ g; \quad
$4$ L $= 400$ cL; \quad $350$ cL $= 3$ L e $50$ cL.
\end{solution}

\begin{solution}{exo:g3:measure:6}
$1$ kg $= 1\,000$ g $> 750$ g: sim, ela tem farinha suficiente, e vão
sobrar $1\,000 - 750 = 250$ g.
\end{solution}

\begin{solution}{exo:g3:measure:7}
Ponteiro pequeno entre $2$ e $3$, ponteiro grande no $6$: $2{:}30$.
Ponteiro pequeno no $9$, grande no $12$: $9{:}00$. Ponteiro pequeno
quase no $5$, grande no $10$: $4{:}50$ (``dez para as cinco'' --- o
ponteiro das horas ainda não chegou ao $5$).
\end{solution}

\begin{solution}{exo:g3:measure:8}
De $8{:}30$ até $9{:}00$: $30$ min; de $9{:}00$ até $11{:}45$:
$2$ h $45$ min. Total: $3$ h $15$ min.
\end{solution}

\begin{solution}{exo:g3:measure:9}
De $9{:}50$ até $10{:}00$: $10$ min; restam $45 - 10 = 35$ min
depois das $10{:}00$. Chegada: $10{:}35$.
\end{solution}

\begin{solution}{exo:g3:measure:10}
$1$ L $= 100$ cL, e $100 = 20 \times 5$: Paulo consegue servir $5$
copos cheios.
\end{solution}

\begin{solution}{exo:g3:measure:11}
Por dia: $600 + 600 = 1\,200$ m. Em cinco dias:
$1\,200 \times 5 = 6\,000$ m $= 6$ km.
\end{solution}
